\section{Задание}

Мы будем работать с данными из датасета\cite{kaggle_russia_real_estate}, в котором содержится информация об объектах недвижимости в России с их стоимостью за 2018--2021 гг.

\begin{itemize}
    \item Цена --- цена в рублях РФ (\textit{числовое});
    \item Дата --- дата публикации объявления (\textit{временное});
    \item Время --- время публикации объявления (\textit{временное});
    \item Координата \(x\) --- значение координаты (широта) (\textit{числовое});
    \item Координата \(y\) --- значение координаты (долгота) (\textit{числовое});
    \item Номер --- код региона РФ (\textit{категориальное});
    \item Тип здания --- тип сооружения, в котором находится квартира (\textit{категориальное});
    \item Этаж --- номер этажа, на котором находится квартира (\textit{категориальное} или \textit{числовое}?);
    \item Число этажей в доме --- общее количество этажей в здании (\textit{категориальное} или \textit{числовое}?);
    \item Комнаты --- количество комнат в квартире, \(-1\) означает студии с открытой планировкой (\textit{категориальное} или \textit{числовое}?);
    \item Площадь --- общая площадь квартиры в кв. метрах (\textit{числовое});
    \item Площадь кухни --- площадь кухни в кв. метрах (\textit{числовое});
    \item Тип квартиры --- тип квартиры, где \(1\) означает вторичный рынок недвижимости, \(11\) --- новостройка (\textit{категориальное}).
\end{itemize}

\textbf{Целевой переменной} будем считать цену (столбец \texttt{price}).

\begin{enumerate}[label=\textbf{Task \arabic*:}, leftmargin=*]

\item \textit{[1 балл]}

Сделайте первичную обработку данных.

\begin{enumerate}[label=\arabic*)]
    \item Выберите один любой регион/город, для которого достаточно данных для построения модели, и оставьте выборку только для этого региона/города. Напишите, какой город выбрали.
    \item Удалите некорректные записи и выбросы.
    \item Создайте дополнительную колонку \texttt{year} --- год продажи.
    \item Удалите колонки \texttt{date} и \texttt{time}.
    \item Исключите \(1\%\) с каждого хвоста распределения по целевой переменной.
    \item Исключите \(1\%\) с каждого хвоста распределения по тем переменным, по которым Вам кажется важно это сделать.
    \item Исследуйте, какие ещё признаки можно преобразовать различными способами, например, нелинейная трансформация и пр. Выберите такие признаки и преобразуйте их.
\end{enumerate}

Выведите на экран первые строки изменённой таблицы, чтобы проверить, что все преобразования сработали.

Исследуйте закон распределения целевого признака. Является ли распределение нормальным?

\item \textit{[1 балл]}

Подготовка к обучению.

\begin{enumerate}[label=\arabic*)]
    \item Создайте матрицу объект-признак \texttt{X} и вектор с целевой переменной \texttt{y}.
    \item Разбейте данные на тренировочную и тестовую части, размер тестовой части --- \(30\%\).
    \item Для линейных моделей важно применить нормализацию, возьмём \texttt{StandardScaler}, обучим метод на тренировочных данных и применим его и к трейну, и к тесту.
    \item После масштабирования признаков \texttt{X\_train} и \texttt{X\_test} перестали быть объектами типа \texttt{pd.DataFrame}, это неудобно. Приведите их снова к \texttt{pd.DataFrame}.
\end{enumerate}

\item \textit{[1 балл]}

Обучите на тренировочных данных с параметрами по умолчанию:
\begin{enumerate}[label=\arabic*)]
    \item линейную регрессию \texttt{Ridge};
    \item метод опорных векторов \texttt{SVR} (с \texttt{kernel='linear'}).
\end{enumerate}

\textbf{Важно!} \texttt{SVR} может долго обучаться без нормализации данных, возьмите на заметку.

С помощью обученных моделей сделайте предсказание на тестовых данных. Вычислите значение метрики \texttt{RMSE} на тесте.

\item \textit{[2 балла]}

Обучите \texttt{DecisionTreeRegressor} с параметрами по умолчанию на тренировочных данных, сделайте предсказание на тесте и вычислите на тесте \texttt{RMSE}.

Нужно ли масштабировать признаки для обучения решающего дерева? Проверьте это, заново разбив данные на тренировочную и тестовую части, назовите новые полученные объекты новыми названиями и проведите обучение, предсказание и оценку метрики до масштабирования и после.

\textit{Подсказка:} можно использовать любой метод масштабирования признаков.

Какой вывод про масштабирование признаков для обучения дерева можно сделать?

\item \textit{[3 балла]}

Теперь ваша цель --- максимально улучшить модели. Для этого:

\begin{enumerate}[label=\arabic*)]
    \item у линейных моделей:
    \begin{itemize}
        \item \texttt{SVR} (если смогли обучить модель) --- подберите гиперпараметр \(C\) (переберите \(20\) значений \(C\) в диапазоне от \(0.1\) до \(10\) с равным шагом);
        \item \texttt{Ridge} --- подберите гиперпараметр \(\alpha\) (переберите \(20\) значений \(\alpha\) в диапазоне от \(10^{-3}\) до \(10^{3}\) с равным логарифмическим шагом).
    \end{itemize}
    \item у решающего дерева подберите \texttt{max\_depth} (переберите глубину от \(1\) до \(10\) с шагом \(1\)) и \texttt{criterion} --- \texttt{squared\_error} или \texttt{absolute\_error}.
\end{enumerate}

Для подбора параметров используйте \texttt{GridSearchCV} с \(5\) фолдами и метрикой \texttt{r2}.

\textbf{Важно!} Подбор параметров осуществляется только по тренировочным данным! Контролируйте свои пайплайны обучения на предмет ликов!

После подбора гиперпараметров обучите модели с найденными гиперпараметрами на тренировочных данных, сделайте предсказания на тесте и оцените их качество (\texttt{RMSE}).

Какая модель справилась с решением задачи лучше всего по метрике \texttt{RMSE}?

\textit{Подсказка:} оформите полученные метрики в \texttt{pd.DataFrame}.

\item \textit{[1 балл]}

Постройте диаграммы (горизонтальные бары) важности признаков для линейной регрессии и дерева (для линейных моделей берите модуль значений весов). Сравните между собой списки из трёх самых важных признаков. Совпадают ли они?

\item \textit{[1 балл]}

Сделайте выводы по самой лучшей модели:
\begin{enumerate}[label=\arabic*)]
    \item Какая модель после подбора гиперпараметров лучше всего решает задачу с точки зрения метрики \texttt{RMSE}?
    \item Достаточное ли качество получившейся модели? Можно ли её использовать в проде?
    \item Как можно ещё улучшить получившуюся модель? Как в случае линейной модели, так и для дерева.
\end{enumerate}

\end{enumerate}